\section{Entering Advent}\label{sec:entering}

Advent begins at dusk on the Saturday before the Fourth Sunday before Christmas, when the Church gathers for Evening Prayer~I of the First Sunday of Advent and opens a new liturgical year. Unlike the secular countdown that floods December with premature Christmas music and commercial haste, the Roman Rite asks Christians to slow down, stand watch, and learn desire again. Before the Church kneels at the manger in Bethlehem, she lifts her eyes to the horizon from which the risen Lord will return in glory.

\subsection{How to Use This Handbook Today}

If you are opening these pages before the First Sunday of Advent, begin with the opening chapters and move week by week through the season, keeping the Evergreen Daily Course (Appendix~A) open beside your Bible and your parish schedule.

If Advent has already begun---perhaps you are picking this up on the Second Sunday, on a busy weekday in mid-December, or on the eve of the Great O~Antiphons---do not try to catch up on missed days. Open directly to the week in which you stand, or to the specific date if December~17 has arrived. The liturgy meets you on the day God has actually given you.

\begin{itemize}
  \item \textbf{Before December~17:} Find your week (First, Second, or Third) and your Sunday Lectionary cycle (Year~A, Year~B, or Year~C), then pair the day's Mass readings with the brief prayer rule in Section~\ref{sec:mass-hours-home} and the matching row of Appendix~A.
  \item \textbf{From December~17 to December~24:} Move to Section~\ref{sec:dec17-24}. Every weekday from December~17 onward has its own date-fixed Scripture readings and, at Evening Prayer from December~17 through December~23, one of the seven Great O~Antiphons.
  \item \textbf{When a major Marian feast arrives:} Turn to Section~\ref{sec:law-calendar} and Section~\ref{sec:week2} for December~8 (the Solemnity of the Immaculate Conception, Patronal Feast of the United States) and December~12 (the Feast of Our Lady of Guadalupe, Patroness of the Americas).
\end{itemize}

\subsection{Why Postconciliar Advent Calls for Its Own Companion}

Readers who know the Roman Rite of 1962 sometimes assume that Advent after the Second Vatican Council is the same season merely translated into English. Opening the liturgical books side by side shows why a dedicated volume for the postconciliar Roman Rite is indispensable:

\begin{enumerate}
  \item \textbf{An explicit two-phase season.} Under the \textit{Universal Norms on the Liturgical Year and the Calendar} (no.~39) and the \textit{Ordo Lectionum Missae}, postconciliar Advent is formally divided into two unequal parts. From Evening Prayer~I of the First Sunday through December~16, the prayers and readings direct the mind toward Christ's glorious coming at the end of time. From December~17 through the morning of December~24, the calendar shifts from week-numbered days to date-numbered privileged weekdays aimed directly at the Nativity, with the Fourth Sunday of Advent always falling inside December~18--24.
  \item \textbf{A three-year Sunday Lectionary.} Where the 1962 Missal repeats the same four Epistles and four Gospels every Advent, the postconciliar \textit{Lectionary for Mass} gives each Sunday a three-reading harmony---Old Testament prophecy, Responsorial Psalm or Canticle, Apostolic Epistle, and Gospel---rotating across Year~A (Matthew), Year~B (Mark, joined by John and Luke), and Year~C (Luke). Over three Decembers, the parish hears twelve distinct Sunday assemblies of Scripture rather than four.
  \item \textbf{A daily weekday Mass course.} Where the 1962 Missal repeats the Sunday Mass on ordinary Advent weekdays (with proper Masses on the three Winter Ember Days and the Vigil of Christmas), the postconciliar Lectionary supplies seventeen week-keyed weekday Scripture pairings for Weeks~I--III (Monday through Saturday in Weeks~I and~II, Monday through Friday in Week~III) plus eight date-keyed pairings for December~17--24, weaving a continuous reading of Isaiah with Advent Gospels before turning on December~17 to the infancy narratives of Matthew and Luke.
  \item \textbf{The revised Liturgy of the Hours.} Promulgated by Pope Paul~VI in \textit{Laudis canticum} (1970), the postconciliar \textit{Liturgy of the Hours} centers lay and parochial prayer on Morning Prayer and Evening Prayer as the hinges of the day, opens the Office of Readings to a richer cycle of Isaiah and patristic homilies, and explicitly preserves the Great O~Antiphons as the evening climax of December~17--23.
  \item \textbf{The reconfiguration of Ember Days and seasonal penance.} The postconciliar General Roman Calendar removed the four universal Ember weeks, transferring the observance of seasonal days of prayer and thanksgiving to episcopal conferences (\textit{Universal Norms}, nos.~45--47). Under the 1983 \textit{Code of Canon Law} (cc.~1250--1253) and United States Conference of Catholic Bishops norms, Advent carries no universal obligation of fasting or abstinence beyond the regular Fridays of the year, shifting the season's ascetical character from statutory fasts to voluntary sobriety, conversion, confession, and almsgiving.
  \item \textbf{The United States Marian calendar.} In the dioceses of the United States, December~8 (the Immaculate Conception) is the nation's Patronal Solemnity and a Holy Day of Obligation (subject to exact rules when December~8 falls on the Second Sunday of Advent and transfers to Monday, December~9), while December~12 is celebrated as the Feast of Our Lady of Guadalupe, Patroness of the Americas---a higher rank than its optional memorial on the General Roman Calendar.
\end{enumerate}

The postconciliar and 1962 calendars share the same ancient root, the same longing for the Lord who comes, and many of the same psalms, antiphons, and Fathers. Yet their lectionaries, weekly architectures, weekday rules, and disciplinary frameworks diverge enough that a faithful Christian keeping the postconciliar Roman Rite needs a handbook built directly from the postconciliar books.
